\section{Introduction}

Ce document présente l'avancement de l'examen machine de programmation parallèle consacré à la simulation d'une galaxie. Le programme considère un trou noir central massif et un grand nombre d'étoiles. À chaque pas de temps, il calcule les accélérations gravitationnelles, puis met à jour les positions et les vitesses au moyen d'un schéma de Verlet avant d'afficher la scène en trois dimensions.

Le travail est organisé par étapes afin de conserver une version exécutable de chaque implémentation intermédiaire. Le mot \og étape \fg{} désigne ici une version précise du projet, avec un objectif technique bien délimité.

\subsection{Signification des étapes}

\begin{itemize}
    \item \textbf{Étape 00 -- baseline :} version de référence, séquentielle, utilisée pour comprendre le code, vérifier le comportement de la simulation et mesurer les temps initiaux.
    \item \textbf{Étape 01 -- Numba avec threads :} version \textit{shared memory}, où l'on parallélise seulement les parties sûres du calcul avec \texttt{Numba} et \texttt{prange}.
    \item \textbf{Étape 02 -- séparation affichage / calcul :} version prévue avec deux processus MPI, un pour l'affichage et un pour le calcul, afin d'éviter que le rendu bloque le calcul.
    \item \textbf{Étape 03 -- MPI complet sur la grille :} version prévue en mémoire distribuée, où plusieurs processus se partagent la grille de calcul avec échange de cellules fantômes.
\end{itemize}

Dans ce rapport, l'étape étudiée en détail est l'étape 01. Les étapes 02 et 03 sont déjà prévues dans l'organisation du dépôt, mais elles seront traitées dans la suite du travail.
